\paragraph{平方根构造。}因子化推导（\cref{sec:tmd-subtraction}）给出的是 $b$ 空间中的乘积 $\tilde f_A^{\mathrm{unsub}}(y_B)\,\tilde f_B^{\mathrm{unsub}}(y_A)/\tilde S(y_A,y_B)$：两个未减除分布在中间快度区重复计算了软辐射，除以软因子正好扣除这一重叠（\cref{fig:tmd-rapidity-bars}）。为了得到两个各自只依赖于一个强子的因子，Collins 把 $\tilde S$ 的平方根分配进两个分布，并在任意的中间快度 $y_n$ 处把软辐射“切开”。下面的引理说明这一构造是恒等式，并给出它的动力学内容。

\begin{lemma}[平方根构造是恒等式 \tagR]\label{lem:tmd-sqrt-identity}
（$R_{\mathrm{def}}$：由算符定义直接推出。）设 $\tilde f_A^{\mathrm{unsub}}(y_B)\equiv\tilde f^{\mathrm{unsub}}_{q/A}(x_A,\mathbf b_T;y_{P_A}-y_B)$ 与 $\tilde f_B^{\mathrm{unsub}}(y_A)\equiv\tilde f^{\mathrm{unsub}}_{\bar q/B}(x_B,\mathbf b_T;y_A-y_{P_B})$ 是\cref{def:tmd-unsub}的未减除关联函数（正规子线 $n_B(y_B)$、$n_A(y_A)$），$\tilde S(y_1,y_2)\equiv\tilde S(\mathbf b_T;y_1,y_2)$ 是\cref{def:tmd-soft}的软因子，平方根按微扰级数理解。
\begin{enumerate}[label=(\alph*),leftmargin=2.2em,itemsep=1pt,topsep=2pt]
\item 对任意 $y_n$，令
\begin{equation}\label{eq:tmd-sqrt-def}
\tilde f_A\equiv\tilde f_A^{\mathrm{unsub}}(y_B)\sqrt{\frac{\tilde S(y_A,y_n)}{\tilde S(y_A,y_B)\,\tilde S(y_n,y_B)}},\qquad
\tilde f_B\equiv\tilde f_B^{\mathrm{unsub}}(y_A)\sqrt{\frac{\tilde S(y_n,y_B)}{\tilde S(y_A,y_B)\,\tilde S(y_A,y_n)}},
\end{equation}
则
\begin{equation}\label{eq:tmd-sqrt-identity}
\tilde f_A\,\tilde f_B=\frac{\tilde f_A^{\mathrm{unsub}}(y_B)\,\tilde f_B^{\mathrm{unsub}}(y_A)}{\tilde S(y_A,y_B)},
\end{equation}
% VF-REF-13
与 $y_n$ 无关；$\tilde S(y_1,y_2)$ 只依赖于 $y_1-y_2$。
\item（蕴含关系，\tagR，$R_{\mathrm{cond}}$，前提\cref{asm:appN-tmd}(T2)）若在 $b$ 空间中 $\ln\tilde S(y_1,y_2)=a_S+(y_1-y_2)\tilde K+\order{\ee^{-(y_1-y_2)}}$、$\ln\tilde f_A^{\mathrm{unsub}}(y_B)=\alpha_A-y_B\tilde K+\order{\ee^{y_B-y_{P_A}}}$（同一斜率 $\tilde K$，$B$ 侧同理），则极限 $y_A\to+\infty$、$y_B\to-\infty$ 存在，$\lim\tilde f_A=\exp(\alpha_A-\frac12a_S-y_n\tilde K)$，从而
\begin{equation}\label{eq:tmd-CS-from-sqrt}
\frac{\partial\ln\tilde f_A}{\partial y_n}=-\tilde K,\qquad\text{即}\qquad\frac{\partial\ln\tilde f_A}{\partial\ln\sqrt{\zeta_A}}=\tilde K,\quad \zeta_A\equiv2(x_AP_A^+)^2\ee^{-2y_n};
\end{equation}
若两个斜率不相等，极限不存在。
\end{enumerate}
（(a) 是算符矩阵元之间的代数关系，对任何外态成立；(b) 的前提是对部分子态的微扰陈述。）
\uses{def:tmd-unsub,def:tmd-soft,asm:appN-tmd}
\end{lemma}

\begin{proof}
(a) 两个平方根之积为 $\sqrt{\frac{\tilde S(y_A,y_n)\tilde S(y_n,y_B)}{\tilde S(y_A,y_B)^2\tilde S(y_n,y_B)\tilde S(y_A,y_n)}}=\frac1{\tilde S(y_A,y_B)}$，$y_n$ 完全消去。$\tilde S$ 只依赖于快度差，是\cref{def:tmd-soft}之后证明的 boost 不变性。

(b) 代入所设的渐近形式，
\[
\ln\tilde f_A=\alpha_A-y_B\tilde K+\tfrac12\big[a_S+(y_A-y_n)\tilde K-a_S-(y_A-y_B)\tilde K-a_S-(y_n-y_B)\tilde K\big]+R=\alpha_A-\tfrac12a_S-y_n\tilde K+R,
\]
其中 $y_A$ 与 $y_B$ 的线性项逐项相消，余项 $R=\order{\ee^{y_B-y_{P_A}}}+\order{\ee^{-(y_A-y_n)}}+\order{\ee^{-(y_A-y_B)}}+\order{\ee^{-(y_n-y_B)}}$ 在 $y_A\to+\infty$、$y_B\to-\infty$ 时趋于零。于是极限存在且等于 $\exp(\alpha_A-\frac12a_S-y_n\tilde K)$，对 $y_n$ 求导得\cref{eq:tmd-CS-from-sqrt}的第一式；由 $\ln\sqrt{\zeta_A}=\ln(\sqrt2x_AP_A^+)-y_n$ 得第二式。若 $\tilde f_A^{\mathrm{unsub}}$ 的斜率是 $\tilde K'\neq\tilde K$，上式中残留 $(\tilde K-\tilde K')y_B$，在 $y_B\to-\infty$ 时发散。% check: s10_subtraction_algebra.py (c)
\end{proof}

\paragraph{地位与缺口（(b) 部分）}(a) 是纯代数，(b) 是以\cref{asm:appN-tmd}(T2) 为前提的严格蕴含；其内容等级就是 (T2) 的等级 \tagA。“$\tilde S$ 只依赖于 $y_1-y_2$”是精确的对称性推论；“$\ln\tilde S$ 与 $\ln\tilde f^{\mathrm{unsub}}$ 渐近线性且斜率相同”是动力学陈述：单圈上它由\cref{lem:tmd-soft-oneloop}与\cref{lem:tmd-unsub-oneloop}显式成立（两者的斜率都是 $\tilde K=-\frac{\as C_F}\pi L_b$，对软因子修正项是 $\order{Y\ee^{-2Y}}$），所有阶上则只有 Collins 的全阶论证\cite[第~10.11、13.8~节]{Collins2011}。斜率相等等价于极限 $\lim\tilde f_A$ 的存在，这正是 Collins--Soper 方程的内容。升级途径：从算符定义出发对所有阶证明 (T2)。

%%% BEGIN fig:tmd-rapidity-bars
\begin{figure}[htbp]
\centering
\begin{tikzpicture}[fd,font=\small,
  rb/.style={draw=black!55,line width=0.5pt,rounded corners=1.5pt},
  rlab/.style={font=\footnotesize,inner sep=1pt}]
  \def\yB{0.6} \def\yPB{2.7} \def\yn{6.0} \def\yPA{9.3} \def\yA{11.4}
  % axis and ticks
  \draw[lcaxis] (-0.2,0) -- (12.4,0) node[right] {$y$};
  \foreach \x/\l in {\yB/y_B,\yPB/y_{P_B},\yn/y_n,\yPA/y_{P_A},\yA/y_A}
    {\draw[line width=0.5pt] (\x,0.08) -- (\x,-0.08) node[below=1pt,rlab] {$\l$};}
  % unsubtracted A: y_B ... y_PA  (overlap part near y_B hatched)
  \fill[blue!12] (\yB,0.35) rectangle (\yPA,0.75);
  \fill[pattern=north east lines,pattern color=blue!45] (\yB,0.35) rectangle (\yPB+1.0,0.75);
  \draw[rb] (\yB,0.35) rectangle (\yPA,0.75);
  \node[rlab,anchor=west] at (\yPA+0.15,0.55) {$\tilde f^{\,\mathrm{unsub}}_A(y_B)$};
  % unsubtracted B: y_PB ... y_A  (overlap near y_A hatched)
  \fill[orange!18] (\yPB,0.95) rectangle (\yA,1.35);
  \fill[pattern=north west lines,pattern color=orange!60] (\yPA-1.0,0.95) rectangle (\yA,1.35);
  \draw[rb] (\yPB,0.95) rectangle (\yA,1.35);
  \node[rlab,anchor=east] at (\yPB-0.15,1.15) {$\tilde f^{\,\mathrm{unsub}}_B(y_A)$};
  % soft factor: y_B ... y_A
  \fill[green!15] (\yB,1.55) rectangle (\yA,1.95);
  \draw[rb] (\yB,1.55) rectangle (\yA,1.95);
  \node[rlab] at ({0.5*(\yB+\yA)},1.75) {$\tilde S(y_A,y_B)$};
  % result: split at y_n
  \fill[orange!18] (\yPB,2.35) rectangle (\yn,2.75);
  \fill[blue!12] (\yn,2.35) rectangle (\yPA,2.75);
  \draw[rb] (\yPB,2.35) rectangle (\yn,2.75);
  \draw[rb] (\yn,2.35) rectangle (\yPA,2.75);
  \node[rlab] at ({0.5*(\yPB+\yn)},2.55) {$\tilde f_B(x_B,b_T;\zeta_B)$};
  \node[rlab] at ({0.5*(\yn+\yPA)},2.55) {$\tilde f_A(x_A,b_T;\zeta_A)$};
  \draw[dashed,gray!60!black,line width=0.5pt] (\yn,0.1) -- (\yn,3.05);
  \node[rlab,anchor=west] at (\yPA+0.15,2.55) {$=\tilde f^{\,\mathrm{unsub}}_A\tilde f^{\,\mathrm{unsub}}_B/\tilde S$};
\end{tikzpicture}
\caption{快度覆盖与平方根构造（\cref{lem:tmd-sqrt-identity}；示意）。横轴为辐射胶子的快度。未减除分布 $\tilde f^{\,\mathrm{unsub}}_A(y_B)$ 包含快度在 $y_B$ 与 $y_{P_A}$ 之间的辐射，$\tilde f^{\,\mathrm{unsub}}_B(y_A)$ 包含 $y_{P_B}$ 与 $y_A$ 之间的辐射；软因子 $\tilde S(y_A,y_B)$ 覆盖 $y_B$ 到 $y_A$。乘积除以 $\tilde S(y_A,y_B)$ 之后每个快度区间恰好计一次：在 $y_{P_B}<y<y_{P_A}$ 中为 $1+1-1=1$，在 $y_B<y<y_{P_B}$ 与 $y_{P_A}<y<y_A$ 中为 $1-1=0$（阴影：Wilson 线把辐射伸入另一强子方向的部分，被软因子抵消）；平方根构造在任意的 $y_n$ 处把结果分成两个因子（最上一行），乘积与 $y_n$ 无关，$y_n$ 的变化由 Collins--Soper 方程描述。}
\label{fig:tmd-rapidity-bars}
\end{figure}
%%% END fig:tmd-rapidity-bars

\begin{definition}[Collins 的 TMD 定义]\label{def:tmd-collins}
强子 $A$ 中夸克的 TMD 分布定义为
\begin{equation}\label{eq:tmd-collins-def}
\tilde f_{q/A}(x,\mathbf b_T;\mu,\zeta_A)\equiv\lim_{y_B\to-\infty}\tilde f^{\mathrm{unsub}}_{q/A}(x,\mathbf b_T;y_{P_A}-y_B)\,\sqrt{\frac{\tilde S(\mathbf b_T;+\infty,y_n)}{\tilde S(\mathbf b_T;+\infty,y_B)\,\tilde S(\mathbf b_T;y_n,y_B)}}\;Z_{\mathrm{UV}},
\qquad\zeta_A\equiv M_A^2x^2\ee^{2(y_{P_A}-y_n)}=2(xP_A^+)^2\ee^{-2y_n},
\end{equation}
% VF-TMD-9
其中 $\tilde S(\mathbf b_T;+\infty,\cdot)$ 表示组合中 $y_A\to+\infty$ 的极限（单个因子发散，比值有限），$y_n$ 是任意的中间快度，$Z_{\mathrm{UV}}$ 是 $\MSbar$ 的紫外重整化因子（包括夸克场的 $Z_2$）。强子 $B$ 的分布交换 $\pm$ 分量与 $A\leftrightarrow B$ 后同样定义，$\zeta_B\equiv2(xP_B^-)^2\ee^{2y_n}$。\textbf{极限的次序是定义的一部分}：先在固定的 $\varepsilon$ 下取 $y_A\to+\infty$、$y_B\to-\infty$，再作 $\MSbar$ 重整化，最后令 $\varepsilon\to0$。
\uses{def:tmd-unsub,def:tmd-soft,lem:tmd-sqrt-identity}
\end{definition}

这一定义就是 Collins 的书中的式~(13.106)\cite[式~(13.106)、(13.107)]{Collins2011}换算到本书约定后的形式；极限次序的规定见\cite[第~10.8.2~节]{Collins2011}。几点说明。(1) 由\cref{lem:tmd-sqrt-identity}，快度发散在未减除分布与软因子之间相消，只有组合 $\tilde f$ 与正规子无关；单个的 $\tilde f^{\mathrm{unsub}}$ 与 $\tilde S$ 都依赖于正规子（例如类时线与类空线给出不同的常数），断言单个 TMD 或软函数与正规子无关是错误的。(2) 由 $M_A^2\ee^{2y_{P_A}}=2(P_A^+)^2$，$\zeta_A$ 的两种写法相同；在 Drell--Yan 截面中 $\zeta_A\zeta_B=(2x_Ax_BP_A^+P_B^-)^2=Q^4$（领头幂次），与 $y_n$ 无关。注意文献中有一处把 $\zeta$ 多写了因子 $2$ 的排印错误（附录~Q 第~2~条），不能照抄；同一作者在后来的工作中用的是不含此因子的形式\cite{Aybat:2011ge}，与这里一致\cite[式~(14.33)]{Collins2011}。(3) 单圈时 $Z_{\mathrm{UV}}=Z_FZ_2=1+aS_\varepsilon[-\frac1{\varepsilon^2}+\frac{l_\zeta-3/2}{\varepsilon}]$，$l_\zeta\equiv\ln(\zeta/\mu^2)$，其推导在\cref{prop:tmd-evolution-oneloop}的证明中；它在两种 $S_\varepsilon$ 约定下相同。% VF-TMD-4
(4) Wilson 线的自相互作用在这一组合中相消\cite[第~13.7.2~节]{Collins2011}，这推广了\cref{lem:tmd-soft-oneloop}(c) 的单圈计算。

\begin{remark}[极限次序]\label{rem:tmd-order-of-limits}
极限次序有两层含义。第一层是\cref{def:tmd-collins}中的规定：在固定 $\varepsilon$ 下先取类光极限。\cref{lem:tmd-unsub-oneloop}(c) 给出了显式的理由：非类光线的顶角图正比于 $(\mu^2/\Lambda_B^2)^\varepsilon$，在 $\Real\varepsilon>0$ 时随 $y_B\to-\infty$ 消失；若先展开 $\varepsilon$，每个顶角图留下 $-\frac{\as C_F}{8\pi}\ln^2\frac{\Lambda_B^2}{\mu^2}$（\cref{eq:tmd-vertex-expanded}），它在类光极限下发散。Collins 指出，若坚持相反的次序，就必须补上一个额外的、纯紫外的重整化因子\cite[式~(10.97)]{Collins2011}。第二层与因子化推导中的近似子有关：在\cref{def:appN-tmd-approx}的快度正规化近似子中，进入硬部分的 $A$ 共线动量投影为 $\pi_+^{(y_B)}(k)=(n_B(y_B)\cdot k)\,n_A$，而 $n_B(y_B)\cdot k=k^+-\ee^{2y_B}k^-$；在固定 $x=k^+/P^+$ 下，这样定义的共线因子与光前矩阵元相差相对量 $\ee^{2y_B}k^-/k^+=\ee^{2(y_B-y_k)}$（$y_k$ 为 $k$ 的快度）。对典型的共线动量（$y_k\approx y_{P_A}$）它是 $\ee^{2(y_B-y_{P_A})}$，很小；但在 $y_k$ 接近 $y_B$ 处它不小，只是在 $y_B\to-\infty$ 时逐点趋于零。在除以软因子之后的组合中把 $y_B\to-\infty$ 移到积分号之内，正是\cref{asm:appN-tmd}(T4) 的内容。% check: s10 (d)
\end{remark}

\begin{remark}[其他方案]\label{rem:tmd-schemes}
快度发散还有几种常用的正规化与组织方式：Ji--Ma--Yuan 用略微离开光锥的规范连线\cite{JiMaYuan2005,Ji:2004xq,Ji:2005nu}；Echevarria--Idilbi--Scimemi 的 $\delta$ 正规子\cite{EchevarriaIdilbiScimemi2012,Echevarria:2012js}；Chiu--Jain--Neill--Rothstein 的快度重整化群\cite{ChiuEtAl2012a,ChiuEtAl2012b}；Becher--Neubert 的“共线反常”\cite{BecherNeubert2011}。Collins 与 Rogers 给出了 CSS 原始表述、Collins 表述与若干 SCET 表述之间的联系，并在 $\order{\as^2}$、$\order{\as^3}$ 上用已知结果核对\cite{Collins:2017oxh}（据其摘要）；这些方案之间的等价性是全阶论证（\tagA），只在低阶被显式验证。快度反常维度（即 $\tilde K$）已经算到四圈\cite{Duhr:2022yyp,Moult:2022xzt}（文献结果）；多部分子软因子快度发散的重整化定理见\cite{Vladimirov:2017ksc}。以 $\delta$ 正规子为例（单圈）：类光 Wilson 线的传播子取 $\frac{\ii}{-k^++\ii\delta^+}$、$\frac{\ii}{k^-+\ii\delta^-}$，未减除分布中的干涉项变为 $\frac{2x(1-x)}{(1-x)^2+r^2}$（$r=\delta^+/p^+$），在 $r\to0$ 时给出 $2x[\frac1{1-x}]_++\delta(1-x)\ln\frac{(p^+)^2}{(\delta^+)^2}$；软因子实胶子的快度积分 $\int\dd y\,\frac{l_T^2}{(l^++\delta^+)(l^-+\delta^-)}=2\ln\frac{l_T^2}{2\delta^+\delta^-}+\order{\delta}$；于是 $\tilde f^{\mathrm{unsub}}/\sqrt{\tilde S}$ 的单圈 $\delta(1-x)$ 项为
\begin{equation}\label{eq:tmd-delta-reg}
\ln\frac{(p^+)^2}{(\delta^+)^2}-\ln\frac{k_T^2}{2\delta^+\delta^-}=\ln\frac{\zeta}{k_T^2},\qquad\zeta=2(p^+)^2\frac{\delta^-}{\delta^+},
\end{equation}
% VF-TMD-8; check: s06 (9)
与 Collins 定义的单圈结果（\cref{eq:tmd-1loop-kspace}）相同，$\zeta$ 与 $\delta^\pm$ 的对应正是 Collins--Rogers 给出的对应\cite{Collins:2012uy}。软函数的欧氏（离光锥）实现与时序无关性见第~23~章。
\end{remark}

\begin{remark}[方案常数（$S_\varepsilon$ 约定）]\label{rem:tmd-scheme-constants}
TMD 与硬因子每圈含 $\varepsilon^{-2}$ 的紫外极点，因此 $\MSbar$ 中 $S_\varepsilon$ 的选取会改变有限常数：$S_\varepsilon/S^{\mathrm C}_\varepsilon=\ee^{-\varepsilon\gamma_E}\Gamma(1-\varepsilon)=1+\frac{\pi^2}{12}\varepsilon^2+\order{\varepsilon^3}$。本书用 $S_\varepsilon=(4\pi)^\varepsilon\ee^{-\varepsilon\gamma_E}$，Collins 用 $S^{\mathrm C}_\varepsilon=(4\pi)^\varepsilon/\Gamma(1-\varepsilon)$\cite{Collins2011}。以 $a=\frac{\as C_F}{2\pi}$ 为单位，两种约定下的常数列于\cref{tab:tmd-scheme-constants}。两者之差在截面中相消：$\mathcal H^{\mathrm{bind}}-\mathcal H^{\mathrm C}=a\frac{\pi^2}{6}$，而 $2[\tilde C^{\mathrm{bind}}_{q/q}-\tilde C^{\mathrm C}_{q/q}]=-a\frac{\pi^2}{6}\delta(1-x)$；$\gamma_q$、$\gamma_K$、$\tilde K$ 在单圈上相同，CSS 系数 $C^{(1)}_{qq}=\frac{\as}{\pi}[\frac{C_F}{2}(1-z)+C_F(\frac{\pi^2}{4}-2)\delta(1-z)]$ 在两种约定下也相同。本书的数值与 Echevarria--Idilbi--Scimemi 的一致，Collins 约定下的数值与 Aybat--Rogers 的一致\cite{EchevarriaIdilbiScimemi2012,AybatRogers2011}；两种硬因子中 $\pi^2$ 项之比 $\frac76$ 正是 Collins 与 Rogers 指出的约定差别\cite{Collins:2012uy}。\textbf{本书的规则}：用本书的 $S_\varepsilon$；引用 Collins 约定下的常数时必须换算并注明。
\end{remark}

\begin{table}[htbp]
\centering
\caption{单圈有限常数对 $\MSbar$ 中 $S_\varepsilon$ 约定的依赖（$a=\as C_F/2\pi$，$L=\ln(Q^2/\mu^2)$，$\tilde C_{q/q}$ 取 $\mu=\mu_b$、$\zeta=\mu_b^2$）。截面中的乘积 $\mathcal H\,\tilde C\,\tilde C$ 与约定无关。}
\label{tab:tmd-scheme-constants}
\small
\begin{tabular}{@{}lcc@{}}
\toprule
 & 本书 $S_\varepsilon=(4\pi)^\varepsilon\ee^{-\varepsilon\gamma_E}$ & Collins $S^{\mathrm C}_\varepsilon=(4\pi)^\varepsilon/\Gamma(1-\varepsilon)$\\
\midrule
$\tilde C_{q/q}(x;\mu_b,\mu_b^2)$ & $\delta(1-x)+a\big[(1-x)-\frac{\pi^2}{12}\delta(1-x)\big]$ & $\delta(1-x)+a(1-x)$\\
$\mathcal H_{\mathrm{DY}}$ & $1+a\big[-L^2+3L-8+\frac{7\pi^2}{6}\big]$ & $1+a\big[-L^2+3L-8+\pi^2\big]$\\
$\mathcal H_{\mathrm{SIDIS}}$ & $1+a\big[-L^2+3L-8+\frac{\pi^2}{6}\big]$ & $1+a\big[-L^2+3L-8\big]$\\
$\gamma_q,\ \gamma_K,\ \tilde K$ & \multicolumn{2}{c}{单圈相同（\cref{prop:tmd-evolution-oneloop}）}\\
\bottomrule
\end{tabular}
\end{table}
% VF-TMD-3, VF-TMD-5; check: s06 (3), (6)
